\section{System Design using State Machine}
\subsection{Overall System State Machine}

The global system is modeled using an asynchronous, event-driven architecture utilizing FreeRTOS synchronization
primitives (semaphores) to achieve non-blocking concurrency.

\begin{itemize}
  \item \textbf{Continuous Monitoring Thread:} Operates as a high-priority periodic task. It updates ambient metrics
        continuously, flashing the status via the HEATER safety profile and refreshing the local LCD. Upon capturing
        stable data, it dispatches an activation token via Release \texttt{sensor\_ready\_sem()} to wake the decision-making
        engine.
  \item \textbf{Dynamic Fork-Join Concurrency Control:} The DECISION MAKER parses the sensor data and evaluates
        constraints. To execute adjustments concurrently without freezing peripheral systems, control flow forks into
        three prospective branches depending on active variables: thermal adjustment (COOLER), moisture adjustment
        (HUMIDIFIER), or an active bypass (Metrics are within the threshold).
  \item \textbf{Adaptive Synchronization Barrier:} Once paths are dispatched, the system blocks at a Join synchronization
        point. The feedback path safely routes control back to the baseline evaluation loop by implementing a dynamic
        conditional check: Acquire necessary semaphores if any. This guarantees that if an actuator runs, the
        decision-maker blocks until its fixed operational window finishes, while simultaneously ensuring that an idle
        loop instantly recycles without stalling.
\end{itemize}

\subsection{Heater Control State Machine}

This state machine manages the system safety levels. It controls the indicator LEDs (D3 and D4) using five distinct
temperature boundaries as shown in the figure.

\textbf{States Description}

\begin{itemize}
  \item \textbf{SAFE GREEN}: Default initial state on power-up, active during optimal ambient temperature conditions.
  \item \textbf{WARNING ORANGE}: Triggered during moderate thermal deviations, indicates active climate regulation is required.
  \item \textbf{CRITICAL RED}: Triggered during extreme thermal conditions, indicates high-risk environments requiring
        immediate intervention.
\end{itemize}

\textbf{Transition Conditions}

\begin{itemize}
  \item \textbf{Initialization}: System boots directly into \textbf{SAFE GREEN}.
  \item \textbf{From SAFE GREEN}: Remains in state on a, moves to \textbf{WARNING ORANGE} on b OR c, moves to
        \textbf{CRITICAL RED} on d OR e.
  \item \textbf{From WARNING ORANGE}: Remains in state on b or c, moves to \textbf{CRITICAL RED} on d OR e, moves to
        \textbf{SAFE GREEN} on a.
  \item \textbf{From CRITICAL RED}: Remains in state on d or e, moves to \textbf{WARNING ORANGE} on b OR c, moves to
        \textbf{SAFE GREEN} on a.
\end{itemize}

\subsection{Cooler Control State Machine}

\textbf{States Description}

\begin{itemize}
  \item \textbf{IDLE BLACK}: Default initial state when the system boots up, the cooler hardware remains completely
        powered \textbf{OFF}.
  \item \textbf{COOLING GREEN}: Active state indicating the physical cooler is running, the cooler runs for a strict,
        non-blocking time window.
\end{itemize}

\textbf{Thresholds and Transition Conditions}

\begin{itemize}
  \item \textbf{Initialization}: The state machine directly enters \textbf{IDLE BLACK} upon system startup.
  \item \textbf{Activation Threshold}: The system transitions from \textbf{IDLE BLACK} to \textbf{COOLING GREEN} if the
        evaluated condition \texttt{[Temp \textgreater{} 28°C]} is met. This immediately triggers the physical
        cooler hardware to turn \textbf{ON}.
  \item \textbf{Fixed Duration}: Upon entering the \textbf{COOLING GREEN} state, a dedicated hardware or software timer
        begins counting. The cooler stays continuously active for exactly \textbf{5 seconds}.
  \item \textbf{Re-check Condition}: Once the condition \texttt{[Timer == 5s]} evaluates to true, the state machine
        turns the cooler \textbf{OFF} and transitions straight back to \textbf{IDLE BLACK}. On the very next clock cycle
        in the IDLE state, the temperature is re-checked against the 28°C activation threshold. If it is still hot, it
        cycles back into cooling; otherwise, it remains safely idle.
\end{itemize}

\subsection{Humidifier Control State Machine}

\textbf{States Description}

\begin{itemize}
  \item \textbf{IDLE BLACK}: Default initial state upon hardware system power-up, the physical humidifier hardware
        remains powered \textbf{OFF}.
  \item \textbf{GREEN}: The initial phase of active humidification running for an exact, timed duration.
  \item \textbf{YELLOW}: The secondary intermediate phase of the active execution cycle.
  \item \textbf{RED}: The final phase of the humidification sequence before monitoring checks reset.
\end{itemize}

\textbf{Activation Threshold and Transition Conditions}

\begin{itemize}
  \item \textbf{Initialization}: The state machine directly initializes into the \textbf{IDLE BLACK} state on startup.
  \item \textbf{Activation Threshold}: The system transitions from \textbf{IDLE BLACK} to \textbf{GREEN} when the
        condition \texttt{[Humi \textless{} 40]} evaluates to true. This transition instantly executes the action
        \textbf{Turn ON Humidifier}.
  \item \textbf{Phase 1 to Phase 2 (GREEN to YELLOW)}: The system remains in the GREEN phase for a fixed duration. Once
        the internal tracking variable meets the condition \texttt{[Time == 5s]}, it transitions to \textbf{YELLOW}.
  \item \textbf{Phase 2 to Phase 3 (YELLOW to RED)}: The system holds in the YELLOW state for an intermediate window.
        Once the condition \texttt{[Time == 3s]} evaluates to true, it transitions straight to \textbf{RED}.
  \item \textbf{Cycle Completion and Re-check (RED to IDLE BLACK)}: The system processes its final active interval in
        the RED state. When the condition \texttt{[Time == 2s]} is reached, the state machine drops back into
        \textbf{IDLE BLACK}.
\end{itemize}

In the immediate next execution clock cycle inside \textbf{IDLE BLACK}, the relative humidity is re-evaluated. If
\texttt{Humi < 40} remains true, the multi-phase cycle repeats immediately starting at \textbf{GREEN}; otherwise, the
hardware drops out of operation and settles in the idle baseline.
